\chapter{Related Works}

\section{Learning Management Platforms}
Traditional LMS services are centered around managing courses, assignment management, and grading. Their operational strength is strong.
but AI-assisted learning is external or less connected. This forms a general disjunction: students are able to.
access information, yet contextual tutoring is not combined with the same course data.

\section{AI-Assisted Learning Systems}
Recent AI tutoring systems demonstrate the importance of interactive question answering, guided explanation, and adaptive practice.
Nonetheless, most of the implementations rely on generic chat interactions with little grounding in course. In educational settings,
examples are more effective when associated with real class content and subject matter.

|human|>section Retrieval-Augmented Generation in Education.
Retrieval-Augmented Generation (RAG) is a combination of language generation with information retrieval. This is in course contexts.
helpful since answers can be linked to existing resources, rather than model memory alone. Current research
promotes hybrid retrieval strategies integrating lexical and vector cues \cite{lewis2020rag,gao2023ragsurvey).
	
	In this project, the corresponding notion is not an originality in designing algorithms, but a practical application of retrieval to.
	Answers to course context in LMS.
	
	OCR in Academic Processes.
	Digitization of documents has been a long term use of OCR. Old-fashioned OCR engines are effective in clean printed text, whereas.
	OCR with methods based on transformers is under investigation in more complicated cases. OCR can be of use in the educational processes.
	to extract text on uploaded content or timetable image to allow content to be indexed and re-used.
section Agentic and Multi-Mode Tutor Interfaces.
In many instances, a modern tutor interface will be a combination of more than one interaction pattern: free chat, explain mode, question answering, quiz.
resource recommendations, practice and. The associated body of literature on tool-using agents and structured prompting.
This direction is informed by cite: react-yao2023 and toolformer-schick2023. In this case, the focus of this report is on implemented.
instead of hypothetical agent performance assertions.

section. Location of this Project.
According to the context above, Cognitive Copilot is a positioned implementation-based academic system that:
\begin{itemize}
	item combines LMS modules and AI tutor processes into a single platform,
	item connects AI interactions to course/topic context,
	item assists with real-world classroom activities (community, notifications, routine, analytics),
	item and emphasizes correctness and maintainability of software.
\end{itemize}

Engineering integration and interoperability of modules is the project contribution but not a claim of a.
new AI algorithm.